\documentclass[11pt,a4paper]{article}
\usepackage[utf8]{inputenc}
\usepackage[T1]{fontenc}
\usepackage{lmodern}
\usepackage[left=2.25cm,right=2.25cm,top=1.95cm,bottom=2.05cm]{geometry}
\usepackage{amsmath,amssymb,booktabs,array,tabularx,graphicx,xcolor,fancyhdr,enumitem}
\usepackage{caption,microtype,eso-pic}
\usepackage[hidelinks]{hyperref}
\renewcommand{\figurename}{Figura}
\renewcommand{\tablename}{Tabla}
\definecolor{azul}{HTML}{233D54}
\definecolor{gris}{HTML}{565E64}
\setlength{\parindent}{0pt}
\setlength{\parskip}{0.43em}
\setlength{\tabcolsep}{4pt}
\renewcommand{\arraystretch}{1.18}
\setlist[itemize]{leftmargin=1.2em,itemsep=2pt,topsep=2pt}
\captionsetup{font=small,labelfont=bf,skip=5pt}
\pagestyle{fancy}
\fancyhf{}
\fancyhead[L]{\small\color{gris}MCDI501 \;|\; Sumativa 2: validación y remuestreo}
\fancyhead[R]{\small\color{gris}Grupo 7}
\fancyfoot[L]{\footnotesize\color{gris}Universidad Andrés Bello \;|\; Magíster en Ciencia de Datos e IA}
\fancyfoot[R]{\footnotesize\thepage\,/\,7}
\setlength{\headheight}{14pt}
\renewcommand{\headrulewidth}{0.3pt}
\newcommand{\titulo}[2]{\vspace*{0.1em}{\Large\bfseries\color{azul}#1}\par{\small\color{gris}#2}\par\vspace{0.4em}}
\newcommand{\subtitulo}[1]{\par\vspace{0.2em}{\bfseries\color{azul}#1}\par}
\newcommand{\figura}[3]{\begin{center}\includegraphics[width=\linewidth]{../../results/03_sumativa2/figures/#1}\captionof{figure}{#2}\label{#3}\end{center}}
\newcommand{\fuente}{\vspace{-0.3em}{\footnotesize\color{gris}Fuente: elaboración del Grupo 7 a partir de \texttt{weatherAUS.csv}.}\par}
% Generado desde resultados_validados_s2.json; no editar cifras a mano.
\newcommand{\PorcentajeLluvia}{22,42}
\newcommand{\FaltantesSol}{48,01}
\newcommand{\DMedia}{20,0769}
\newcommand{\DSd}{5,5889}
\newcommand{\DIntervaloT}{[18,472; 21,682]}
\newcommand{\TValor}{25,1461}
\newcommand{\PValorSignos}{0,00009999}
\newcommand{\AsimetriaD}{-0,0757}
\newcommand{\ShapiroW}{0,9884}
\newcommand{\ShapiroP}{0,9059}
\newcommand{\MediaMC}{20,07287}
\newcommand{\SdMC}{0,79417}
\newcommand{\SdTeorica}{0,79841}
\newcommand{\IntervaloMC}{[18,503; 21,618]}
\newcommand{\ErrorMC}{0,00794}
\newcommand{\DiferenciaMC}{-0,00403}
\newcommand{\ErrorUnidadesMC}{-0,51}
\newcommand{\DIntervaloBootstrap}{[18,517; 21,617]}
\newcommand{\JackMin}{19,8007}
\newcommand{\JackMax}{20,3584}
\newcommand{\JackAlbany}{+0,2815}
\newcommand{\JackAlice}{-0,2762}
\newcommand{\JackCobar}{-0,1961}
\newcommand{\JackCambioMax}{0,2815}
\newcommand{\JackCambioRel}{1,40}
\newcommand{\RhoLluvia}{0,4436}
\newcommand{\RhoTemp}{0,9845}
\newcommand{\RhoSol}{-0,6203}
\newcommand{\RhoLluviaFiltro}{0,3074}
\newcommand{\CambioRhoLluvia}{-0,1362}
\newcommand{\ReduccionRhoLluvia}{30,71}
\newcommand{\ExcluidosLluvia}{24.790}
\newcommand{\PorcentajeExcluidoLluvia}{17,89}
\newcommand{\CercaLluvia}{2,00}
\newcommand{\PValorT}{2{,}7471\times10^{-29}}
\newcommand{\LagMediana}{0,5023}
\newcommand{\MediaMCmil}{20,05495}
\newcommand{\MediaMCcincomil}{20,07428}
\newcommand{\CambioPercentilHum}{-2,17}
\newcommand{\CambioBcaHum}{-3,56}
\newcommand{\CambioPercentilMax}{-3,91}
\newcommand{\CambioBcaMax}{-2,50}
\newcommand{\CambioPercentilMin}{-3,61}
\newcommand{\CambioBcaMin}{-2,21}
\newcommand{\DesplazaBcaHum}{-0,213}
\newcommand{\DesplazaBcaMin}{+0,136}
\newcommand{\AnchoCorrTemp}{0,0056}
\newcommand{\AnchoCorrSol}{0,0792}
\newcommand{\AnchoCorrLluvia}{0,0632}
\newcommand{\TablaBootstrap}{\begin{tabular}{@{}lrrrr@{}}\toprule Variable & Media S1 & IC t & IC percentil & IC BCa\\\midrule
Humidity3pm (\%) & 51,247 & [48,025; 54,469] & [48,045; 54,349] & [48,042; 54,256]\\
MaxTemp ($^{\circ}$C) & 23,388 & [22,223; 24,553] & [22,275; 24,514] & [22,300; 24,571]\\
MinTemp ($^{\circ}$C) & 12,259 & [11,124; 13,394] & [11,188; 13,376] & [11,260; 13,480]\\
\bottomrule\end{tabular}}
\newcommand{\TablaCorrelaciones}{\begin{tabular}{@{}lrrrr@{}}\toprule Par & $n$ pares & Localidades & $\rho$ S1 & IC bootstrap\\\midrule
MaxTemp / Temp3pm & 141.589 & 49 & 0,9845 & [0,9812; 0,9868]\\
Humedad / Sunshine & 74.497 & 30 & -0,6203 & [-0,6595; -0,5803]\\
Humedad / Rainfall & 138.569 & 49 & 0,4436 & [0,4107; 0,4739]\\
\bottomrule\end{tabular}}
\newcommand{\TablaRobustez}{\begin{tabular}{@{}lrrrr@{}}\toprule Par & Spearman & Kendall $\tau_b$ & Spearman filtro & Excluidos\\\midrule
MaxTemp / Temp3pm & 0,9845 & 0,9117 & 0,9842 & 0,55\,\%\\
Humedad / Sunshine & -0,6203 & -0,4492 & -0,6203 & 0,00\,\%\\
Humedad / Rainfall & 0,4436 & 0,3403 & 0,3074 & 17,89\,\%\\
\bottomrule\end{tabular}}

\begin{document}

% Portada editable: estructura y posiciones del PDF de Sumativa 1 del Grupo 7.
% Las coordenadas están en bp, medidas desde la esquina superior izquierda.
% El logotipo es el recurso original extraído sin alteraciones de ese PDF.
\newcommand{\portadaEvaluacion}{Evaluación Sumativa 2}
\newcommand{\portadaSubtitulo}{Validación, Simulación y Métodos de Remuestreo}
\newcommand{\portadaProyectoUno}{Relación entre condiciones meteorológicas y lluvia al día}
\newcommand{\portadaProyectoDos}{siguiente}
\newcommand{\portadaFecha}{18 de octubre de 2026}
\newcommand{\portadaRepo}{https://github.com/BenjaArayaM/Proyecto-Lluvias-Australia-Grupo7}
\newcommand{\portadaTexto}[3]{%
  \put(297.66,-#1){\makebox[0pt][c]{%
    \fontfamily{ptm}\fontsize{#2bp}{#2bp}\selectfont\color{black}#3}}}
\AddToShipoutPictureFG*{\AtPageUpperLeft{%
  \setlength{\unitlength}{1bp}%
  \put(297.60,-178.80){\makebox[0pt][c]{%
    \includegraphics[width=125.28bp,height=103.92bp]{logo_unab.jpeg}}}%
  \portadaTexto{203.40}{12}{UNIVERSIDAD ANDRÉS BELLO}%
  \portadaTexto{223.32}{11.52}{Magíster en Ciencia de Datos e Inteligencia Artificial}%
  \portadaTexto{251.28}{11.52}{MCDI501}%
  \portadaTexto{265.08}{11.52}{Estadística Computacional para la Toma de Decisiones}%
  \portadaTexto{309.36}{15}{\bfseries\portadaEvaluacion}%
  \portadaTexto{332.64}{12}{\portadaSubtitulo}%
  \portadaTexto{369.72}{18}{\portadaProyectoUno}%
  \portadaTexto{390.48}{18}{\portadaProyectoDos}%
  \portadaTexto{419.64}{11.04}{Informe técnico grupal · Grupo 7 · Ponderación 20 \%}%
  \portadaTexto{451.92}{11.04}{\bfseries Integrantes}%
  \portadaTexto{470.16}{11.04}{Benjamín Araya}%
  \portadaTexto{483.48}{11.04}{Alan Cañete}%
  \portadaTexto{496.80}{11.04}{Patricio Cortez}%
  \portadaTexto{510.00}{11.04}{Jorge Guaico}%
  \portadaTexto{541.32}{11.04}{Dataset: Rain in Australia --- Weather Australia}%
  \portadaTexto{554.64}{11.04}{Docente: Jean Paul Maidana}%
  \portadaTexto{567.84}{11.04}{Fecha: \portadaFecha}%
  \portadaTexto{594.12}{11.04}{\bfseries Repositorio y rama del análisis}%
  \portadaTexto{612.48}{11.04}{%
    \href{\portadaRepo}{\color{blue}\underline{\portadaRepo}}}%
}}
\thispagestyle{empty}
\null
\clearpage


\titulo{1. Bootstrap de tres parámetros}{IC individuales del 95\,\% y comparación con la Sumativa 1}
Se validan resultados propios de S1 (Araya et al., 2026), según jpmaidana (2026).
Rain in Australia (jsphyg, s.\,f.) contiene 145.460 filas, 23 columnas y 49 localidades
(2007-11-01 a 2017-06-25). Se conserva el CSV sin imputación ni eliminación global.

\textbf{Procedimiento.} El estimando $\widehat\theta_X=49^{-1}\sum_\ell\overline X_\ell$
da igual peso a las medias históricas locales. Se remuestrean 49 medias con reemplazo,
10.000 veces por variable, y se calculan IC percentil y BCa (SciPy community, s.\,f.-a).
Se comparan con $\widehat\theta\pm t_{0,975;48}s/\sqrt{49}$ de S1.

\begin{center}
\captionof{table}{Medias con igual peso por localidad e intervalos del 95\,\%.}
\label{tab:bootstrap}
\small\TablaBootstrap
\end{center}
\fuente
\figura{s2_bootstrap_medias.png}{Distribuciones de las medias bootstrap y límites de los tres métodos. Cada panel utiliza su unidad y escala.}{fig:bootstrap}

\subtitulo{Discrepancias cuantificadas}
El ancho percentil cambia \CambioPercentilHum{}\,\% en humedad,
\CambioPercentilMax{}\,\% en temperatura máxima y \CambioPercentilMin{}\,\% en mínima
frente al IC t. Los cambios BCa son \CambioBcaHum{}\,\%, \CambioBcaMax{}\,\% y
\CambioBcaMin{}\,\%. El extremo superior BCa de humedad se desplaza
\DesplazaBcaHum{} pp; el inferior de temperatura mínima, \DesplazaBcaMin{}\,$^{\circ}$C.
La tabla~\ref{tab:bootstrap} y la figura~\ref{fig:bootstrap} muestran concordancia
de magnitudes e intervalos ampliamente solapados, dentro del modelo empleado.

\subtitulo{Confiabilidad e interpretación}
Agregar por localidad evita tratar los días como réplicas independientes. Se suponen
localidades comparables e independencia aproximada entre ellas, sin diseño
probabilístico nacional ni independencia espacial demostrada. Se condiciona a
los historiales observados: no se simulan nuevos periodos ni se corrige sesgo por faltantes.

BCa ajusta sesgo y aceleración, pero no subsana selección de estaciones o dependencia.
Un intervalo más corto no demuestra mejor cobertura. Los IC son individuales, sin
garantía de cobertura conjunta; no atribuyen probabilidad al parámetro fijo ni
describen variación entre días. Las semillas percentiles son $20261018+i$ y las BCa
$20261118+i$, con $i=0,1,2$. Al usar remuestras distintas, parte de las discrepancias
entre extremos corresponde a variación Monte Carlo.

\clearpage
\titulo{2. Validación mediante cambios de signo}{Prueba bilateral de la diferencia de humedad por localidad}
Se usan 138.583 registros válidos para construir 49 diferencias
$D_\ell=\overline H_{\ell,Yes}-\overline H_{\ell,No}$ de Humidity3pm.
Se retoma $H_0:E(D_\ell)=0$ frente a $H_1:E(D_\ell)\ne0$, con $\alpha=0,05$.
Se genera $T_b^*=49^{-1}\sum_\ell s_{b\ell}D_\ell$, donde los signos son
independientes y equiprobables en $\{-1,+1\}$. El estadístico observado es
$T=\DMedia{}$ pp. Este procedimiento corresponde al esquema de permutación
pareada por cambios de signo descrito por SciPy community (s.\,f.-b).

\begin{center}
\captionof{table}{Comparación de la prueba de S1 y el procedimiento computacional.}
\label{tab:permutacion}
\begin{tabularx}{\linewidth}{@{}Xr@{}}\toprule
Estadístico o resultado & Valor\\\midrule
Localidades; diferencias positivas & 49; 49 de 49\\
Diferencia media e IC t 95\,\% (pp) & \DMedia{}; \DIntervaloT\\
Prueba t; grados de libertad & \TValor{}; 48\\
Valor $p$ paramétrico & $\PValorT$\\
Cambios de signo; réplicas tan extremas & 10.000; 0\\
Valor $p$ Monte Carlo, $(r+1)/(B+1)$ & \PValorSignos\\\bottomrule
\end{tabularx}
\end{center}
\figura{s2_permutacion_diagnostico.png}{Distribución nula simulada y diagnóstico Q-Q de las diferencias centradas. La forma observada no acredita los supuestos de la distribución nula.}{fig:permutacion}

\subtitulo{Supuesto que hace válida la comparación}
La media nula no basta: bajo $H_0$ se requiere invariancia frente a invertir signos,
por ejemplo diferencias independientes y simétricas alrededor de cero. La lluvia
no fue asignada al azar; el intercambio de los grupos no está garantizado por diseño.
Por ello, el resultado se interpreta como contraste condicional a ese modelo y no
como una prueba exacta de aleatorización del estudio observacional.

La asimetría muestral es \AsimetriaD{}; Shapiro--Wilk da $W=\ShapiroW{}$ y
$p=\ShapiroP{}$, sin diferencias locales fuera de las cercas de Tukey. Estos
diagnósticos son compatibles con la aproximación de forma de S1, pero no demuestran
simetría bajo la hipótesis nula ni independencia espacial.

\subtitulo{Concordancia y relevancia práctica}
Ambos enfoques rechazan $H_0$ al 5\,\% bajo sus supuestos. El valor computacional
alcanza el mínimo resoluble con 10.000 réplicas y corrección $+1$; no equivale al
valor paramétrico ni debe reportarse como cero. Se conserva t como referencia de S1
y los cambios de signo como sensibilidad a la forma; ninguno resuelve la dependencia
espacial. No se justifica preferir un método solo por su valor $p$.
La humedad es mayor antes de días con lluvia en las 49 localidades observadas;
\DMedia{} pp de humedad no son un aumento equivalente en probabilidad de lluvia.

\clearpage
\titulo{3. Estabilidad de correlaciones}{Bootstrap de conglomerados completos y límites de generalización}
Se seleccionan tres relaciones de la matriz propia de S1. Para cada par se conservan
casos observados y se extraen localidades con reemplazo; se concatenan sus registros
y se recalculan los rangos de Spearman en cada una de 10.000 réplicas válidas.
Los IC se obtienen por percentiles 2,5 y 97,5\,\%.

\begin{center}
\captionof{table}{Correlaciones agrupadas de S1 y sus IC bootstrap del 95\,\%.}
\label{tab:correlaciones}
\small\TablaCorrelaciones
\end{center}
\fuente
\figura{s2_correlaciones_distribuciones.png}{Distribuciones bootstrap con escalas propias. Las líneas señalan el valor de S1 y los extremos del IC; la inclusión de cero se evalúa en la tabla.}{fig:correlaciones}

\subtitulo{Resultados estables y resultados condicionados}
Los tres IC excluyen cero. La dirección positiva de temperatura máxima--temperatura
de las 15:00 es muy estable, con amplitud \AnchoCorrTemp{}; la asociación próxima
a uno sugiere redundancia que deberá comprobarse en el entrenamiento mediante
diagnóstico de multicolinealidad y comparación de modelos, sin decidirla solo por correlación.

Sunshine tiene \FaltantesSol{}\,\% de faltantes. Humedad--sol conserva signo negativo,
pero usa 74.497 pares en 30 localidades.
Su intervalo, de amplitud \AnchoCorrSol{}, no demuestra que la relación sea igual
en las 19 localidades sin Sunshine observado. Humedad--precipitación conserva signo
positivo, con amplitud \AnchoCorrLluvia{}; su magnitud responde al filtro de rango
de la sección~5. Se distingue estabilidad del signo bajo remuestreo de robustez
general frente a cambios en composición, periodos o extremos.

\subtitulo{Qué dependencia conserva este procedimiento}
Los días de una localidad permanecen juntos cuando se remuestrea su conglomerado;
por tanto, no se los remuestrea como unidades independientes. El procedimiento
condiciona a sus historiales observados, sin generar trayectorias futuras ni corregir
dependencia entre localidades. En el coeficiente agrupado, las estaciones con más
registros aportan más filas: este estimando difiere de una media de correlaciones locales.

La exclusión por pares evita una pérdida global de observaciones, pero cambia
denominadores y composición entre relaciones. Estacionalidad, geografía y patrones
de faltantes pueden explicar parte de las asociaciones agrupadas. Ninguno de estos
coeficientes establece causalidad ni acredita la utilidad predictiva de una variable.

\clearpage
\titulo{4. Simulación Monte Carlo}{Incertidumbre de una media en localidades hipotéticamente comparables}
\textbf{Escenario.} Se estudia cuánto variaría la diferencia media de humedad si se
observaran muestras de 49 localidades comparables. Se usa exclusivamente la media
y desviación de las diferencias de S1:
\[
 D^*_{b\ell}\overset{ind}{\sim}N(\widehat\mu_D,\widehat\sigma_D^2),
 \quad (\widehat\mu_D,\widehat\sigma_D)=(\DMedia{};\,\DSd{})\ \mathrm{pp},
 \quad \overline D_b^*=\frac{1}{49}\sum_{\ell=1}^{49}D^*_{b\ell}.
\]
Se generan 10.000 medias. La normalidad aproxima la distribución entre localidades
de S1; no modela la serie meteorológica diaria. Es un escenario paramétrico condicionado
a estimaciones fijas, sin incorporar incertidumbre adicional en esos parámetros.

\begin{center}
\captionof{table}{Resultados de simulación y referencias analíticas.}
\label{tab:mc}
\begin{tabularx}{\linewidth}{@{}Xr@{}}\toprule
Estadístico & Resultado (pp)\\\midrule
Media simulada; objetivo S1 & \MediaMC{}; \DMedia{}\\
Desviación de medias simuladas & \SdMC\\
Desviación teórica $\widehat\sigma_D/\sqrt{49}$ & \SdTeorica\\
Percentiles 2,5 y 97,5\,\% de medias simuladas & \IntervaloMC\\
Error Monte Carlo de la media, $s(\overline D^*)/\sqrt{10.000}$ & \ErrorMC\\
Media simulada menos objetivo S1 & \DiferenciaMC\\\bottomrule
\end{tabularx}
\end{center}
\figura{s2_montecarlo.png}{Convergencia desde la iteración 100 y distribución de medias hipotéticas. La banda de la izquierda es una referencia puntual del error Monte Carlo, no una banda simultánea.}{fig:mc}

\subtitulo{Convergencia numérica}
La discrepancia final corresponde a \ErrorUnidadesMC{} errores Monte Carlo, menor
que uno en valor absoluto. La dispersión simulada es cercana a la referencia analítica.
Las medias acumuladas con 1.000, 5.000 y 10.000 iteraciones son
\MediaMCmil{}, \MediaMCcincomil{} y \MediaMC{} pp. Los percentiles se estabilizan
aproximadamente, aunque conservan variación de simulación; 10.000 iteraciones no
son por sí solas una demostración de convergencia o adecuación del escenario.

\subtitulo{Interpretación para la decisión}
El intervalo central de medias hipotéticas queda sobre cero: bajo esta aproximación,
la diferencia descriptiva persistiría en localidades comparables. Esto respalda estudiar
la humedad como candidato en una fase posterior. El rango simulado \textbf{no es un IC
del parámetro desconocido}, una probabilidad de lluvia ni una validación del pronóstico.
Su coincidencia con la evidencia anterior depende del modelo normal, la independencia
y la calibración con S1; no constituye información meteorológica independiente.

\clearpage
\titulo{5. Robustez e influencia}{Separar estabilidad del signo de sensibilidad de la magnitud}
\textbf{Diferencias locales.} El bootstrap percentil de 10.000 medias da un IC de
\DIntervaloBootstrap{} pp. Al retirar una localidad cada vez, las medias jackknife
varían entre \JackMin{} y \JackMax{} pp. Albany cambia la media en \JackAlbany{}
pp; AliceSprings, \JackAlice{} pp; Cobar, \JackCobar{} pp. El máximo cambio absoluto
es \JackCambioMax{} pp (\JackCambioRel{}\,\% de la media completa). Son unidades
locales influyentes, no errores demostrados; ninguna exclusión individual cambia el signo.

\begin{center}
\captionof{table}{Correlaciones y sensibilidad al filtro de 1,5 IQR en ambas variables.}
\label{tab:robustez}
\small\TablaRobustez
\end{center}
\fuente
\figura{s2_robustez.png}{Influencia de localidades sobre la media de diferencias y comparación de Spearman con datos completos y restricción diagnóstica.}{fig:robustez}

\subtitulo{Humedad--precipitación: magnitud sensible}
Spearman pasa de \RhoLluvia{} a \RhoLluviaFiltro{}: cambio \CambioRhoLluvia{}
y reducción relativa de magnitud de \ReduccionRhoLluvia{}\,\%. Se excluyen
\ExcluidosLluvia{} de 138.569 pares (\PorcentajeExcluidoLluvia{}\,\%). La cerca
superior de Rainfall es \CercaLluvia{} mm: también retira días de lluvia plausibles,
por lo que modifica el rango y la composición de la muestra. El descenso no demuestra
datos erróneos ni un efecto causal de los extremos. Se concluye \textbf{signo positivo
estable bajo bootstrap y magnitud sensible a restringir la precipitación}.

\subtitulo{Otras relaciones y recomendaciones}
Temperatura máxima--temperatura de las 15:00 cambia muy poco al aplicar el filtro.
En humedad--sol no se retiran filas; la igualdad del coeficiente no aporta por sí sola
una prueba de resistencia a extremos. Kendall tau-b coincide en dirección con Spearman,
pero sus magnitudes no son directamente equivalentes. Se conserva el análisis principal
con valores plausibles; cualquier transformación futura se decidirá dentro del entrenamiento.

El jackknife evalúa influencia individual entre localidades y el filtro evalúa una
restricción de rango. No prueban estabilidad ante todos los climas, faltantes o futuros
periodos. La independencia espacial y la invariancia de signos siguen siendo supuestos
no demostrados; deberán considerarse al diseñar la validación predictiva.

\clearpage
\titulo{6. Síntesis validada para Sumativa 3}{Evidencia conservada y decisiones que deberán comprobarse}
\begin{center}
\captionof{table}{Resultados reutilizables y acciones posteriores. Los IC y coeficientes completos están en las tablas anteriores y el archivo de resultados.}
\label{tab:sintesis}
\small
\begin{tabularx}{\linewidth}{@{}p{3.1cm}XX@{}}\toprule
Resultado & Evidencia y límite & Acción para S3\\\midrule
Tres medias locales & IC t, percentil y BCa concordantes; inferencia condicional. & Conservar estimando y unidades; no imputar con medias calculadas en toda la base.\\
Diferencia de humedad & \DMedia{} pp; IC bootstrap \DIntervaloBootstrap. Influencia máxima: Albany, \JackCambioMax{} pp. & Evaluar humedad disponible a la hora de emisión; comprobar estabilidad temporal y geográfica.\\
MaxTemp / Temp3pm & $\rho=\RhoTemp$; fuerte asociación y escasa sensibilidad al filtro. & Revisar redundancia y disponibilidad horaria antes de combinar predictores.\\
Humedad / Sunshine & $\rho=\RhoSol$; IC negativo; cobertura de 30 localidades. & Comparar estrategias con y sin Sunshine, con idéntica partición.\\
Humedad / Rainfall & $\rho=\RhoLluvia$; filtro: \RhoLluviaFiltro. Signo estable, magnitud sensible. & Conservar precipitación plausible; justificar transformaciones dentro del entrenamiento.\\
Calidad y dependencia & Sunshine ausente en 19 localidades; lluvia \PorcentajeLluvia{}\,\%; lag-1 \LagMediana. & Partición temporal, preprocesamiento dentro de entrenamiento y métricas por clase.\\\bottomrule
\end{tabularx}
\end{center}

\subtitulo{Recomendaciones operativas para la fase siguiente}
\begin{itemize}
\item Fijar hora y horizonte del pronóstico; comprobar qué medidas están disponibles entonces.
\item Separar por fechas completas, evitando que un mismo día se distribuya entre entrenamiento y prueba; evaluar localidades si se busca transferir el modelo.
\item Ajustar imputación, escalado, selección y umbral solo con entrenamiento y su validación. El archivo completo ya fue explorado en fases previas; se declarará ese antecedente.
\item Informar matriz de confusión, precision, recall, F1 y ROC-AUC según la rúbrica de modelamiento; PR-AUC puede complementar. No optimizar decisiones con la prueba final.
\end{itemize}

\textbf{Trazabilidad y continuidad.} Entre 142.193 etiquetas conocidas, 31.877 son
\texttt{Yes}. El notebook concilia S1, conserva semillas y verifica el CSV por SHA-256
(prefijo \texttt{573fd715cd69fcac}). El código base es de Patricio, commit \texttt{1484dd0}.
\path{results/03_sumativa2/resultados_validados_s2.json} y su tabla de síntesis guardan
IC, supuestos e influencia; \path{docs/continuidad_sumativa2.md} registra decisiones.
La evidencia diferencia Work/Windows. Estos insumos preparan S3; no la desarrollan.

\subtitulo{Referencias}
{\footnotesize\setlength{\parskip}{0.32em}
Araya, B., Cañete, A., Cortez, P., \& Guaico, J. (2026). \textit{Sumativa 1: informe técnico de avance y notebook} [Informe académico y código]. GitHub.
\url{https://github.com/BenjaArayaM/Proyecto-Lluvias-Australia-Grupo7/tree/134859256a519e0ffde2827319d73a603598e0ea}.

jpmaidana. (2026). \textit{MCDI501: plantilla Sumativa 2, validación, simulación y remuestreo} [Material docente]. GitHub.
\url{https://github.com/jpmaidana/MCDI501-202682}.

jsphyg. (s.\,f.). \textit{Rain in Australia} [Conjunto de datos]. Kaggle.
\url{https://www.kaggle.com/datasets/jsphyg/weather-dataset-rattle-package}.

SciPy community. (s.\,f.-a). \textit{scipy.stats.bootstrap}. SciPy reference guide.
\url{https://docs.scipy.org/doc/scipy/reference/generated/scipy.stats.bootstrap.html}.

SciPy community. (s.\,f.-b). \textit{scipy.stats.permutation\_test}. SciPy reference guide.
\url{https://docs.scipy.org/doc/scipy/reference/generated/scipy.stats.permutation_test.html}.
}
\end{document}
